\sectionguard
\section*{Source-Grounded Synthesis Across the Propers}

Four claims that stand only when the whole formulary is taken together, and
that no element of it makes on its own: how the Mass's three strands reached
this Sunday, which of its chants were made for it, what the book's general
rules rather than this day's own decision settle about it, and how far the
published English reaches. Guéranger's chapter for this Sunday is devotional
and doctrinal throughout and carries no liturgical history; Schuster carries the
chant history and does not read the Mass as one action.

\subsection*{C1. The Sunday's number is the least stable thing about it, and its three strands were assembled on different schedules}

\textbf{The orations did not arrive as a set.} The Secret and the Postcommunion
are a pair in the older books before either is this Sunday's, and the Collect
reaches them only where a Gregorian book sets all three at one Sunday, so
\textbf{the trio as the 1962 book has it is a composition of the Gregorian
tradition and not an inheritance entire from the Gelasian}. \textbf{Two editors
more than a century apart agree on that from opposite directions}: Gerbert
divides the three by descent, marking the Collect apart from the other two,
where Wilson prints them undivided as one Mass.

\textbf{The offset is not one number, and it moves inside one book.} The
Gregorian stands one Sunday out. The Frankish Gelasian stands two out at the
Fifteenth Sunday and four out here --- \textbf{and Gerbert states the arithmetic
himself} in a footnote to that very page: \latin{quae est hic Dom. XX. in
Missali est XVI. \& apud MURAT. XVII.} The two places that fall between are
printed in his book. One is the September Ember Sunday at his p.~178, headed
\latin{Die Dom. Ad S. PETRUM}, \textbf{whose collect \latin{Absolve, quaesumus
Dne, tuorum delicta populorum} is the collect the 1962 Missal prints at the
Twenty-third Sunday after Pentecost, word for word apart from one word
Gerbert's text carries and the Missal's does not} --- \latin{a peccatorum
nostrorum nexibus} there against \latin{a peccatorum nexibus} in the 1962 book
--- so that Mass did not vanish from the Roman book but moved seven Sundays
later. The other is at his p.~180 under the running head \latin{DOM. XIX. POST PENT.}, headed
\latin{Dominica vacat.}: a Mass with its own items, printed as a Mass, with
Gerbert's footnote that it is missing from Muratori and from the Roman Missal
alike. Wilson's own footnote adds that Pamelius joins the second collect of that
twentieth-Sunday Mass to the same Secret and Postcommunion at his seventeenth
Sunday, which is the number Wilson's Gregorian gives.

\textbf{The chants travel on a different number again.} The antiphonary cues of
Cod.~Ottobonianus lat.~313, in the margins of the Gregorian as Wilson prints
them, put this formulary's Introit, Gradual, Offertory and Communion together
against the heading \latin{Dominica XV post Pentecosten} --- one below the 1962
heading and two below the sacramentary heading carrying the same Sunday's
orations --- and the alleluia they give there is \latin{Laudate dominum}. The
same alignment holds through five consecutive Masses in that manuscript.
\textbf{Two readings of Wilson's footnote reference marks place the cue block
differently}: one puts the set two Sundays further on, another puts the offset at
zero. What does not turn on the attachment is that the four chants travel
together in that manuscript and that its alleluia is not the 1962 book's.

\textbf{The readings do not travel as a set at all.} In the \work{Liber Comitis}
as Migne prints it at PL~30 col.~522, this
Gospel stands under \latin{Dominica mensis viii} --- paired with Ephesians
4:1--6, not with Eph.~3:13--21, which is absent from that lectionary
altogether. \latin{Deficiatis} stands once in the whole text of PL~30,
in a Pelagian commentary among the spuria; the eight Ephesians pericopes that
Comes appoints are Eph.~1:3, 2:19, 4:1, 4:23, 4:29, 5:1, 5:11 and 5:15; a
second independent digitisation reads the same; and neither of the two lacunae
Migne marks falls where this Epistle would stand. \textbf{So at this Sunday the
lectionary is the least stable of the three strands}: the chants and the
orations travel as sets and differ only in the number a book gives them, and
the readings do not travel together at all.

\textbf{A medieval use keeps the split whole, and the English books received
it.} Dickinson's \work{Missale ad usum \dots\ Sarum} puts this formulary's
chants and Epistle at \latin{Dominica Decima sexta post Trinitatem}, joined to
the orations and Gospel of the Roman Fifteenth Sunday, and its Collect and
Gospel at \latin{Dominica Decima septima}, joined to the chants and Epistle of
the Roman Seventeenth: the chant strand one Sunday earlier than the oration
strand, the same direction as the Ottobonianus margin. Its alleluias at those
two Sundays are \latin{Qui timent Dominum sperent in eum} and \latin{Dextera
Domini fecit virtutem}, neither of them the Roman book's --- both incipits, and
the arrangement they belong to, from an uncorrected optical layer of Dickinson's
columns, with no page image of them opened. \textbf{Brightman prints that Sarum
Mass of Trinity~XVII against the English books of 1549, 1552 and 1661 and the
arrangement is unchanged through all three} --- ``Lorde, we
praye thee that thy grace maye alwayes preuent and folowe vs, and make vs
continually to be geuen to all good workes,'' with Eph.~iv and Luke~xiv --- and
Blunt tabulates this formulary's own Epistle, Eph.~iii.~13--21, at the
\emph{Sixteenth} Sunday after Trinity with a different Gospel. \textbf{The
Prayer Book did not create the split; it inherited it.} \textbf{The four
witnesses that agree the Collect is Gregorian do not agree where in the
Gregorian it stands}: Wilson's Hadrianum prints it at the seventeenth Sunday
and Gerbert marks it Gregorian as Muratori has it, at that same number, while
Brightman tags it \latin{(Greg. 172)} and Blunt
\latin{Greg. Orationes Quotidianae}, the Gregorian's daily orations rather than
any Sunday.

\textbf{And a Western rite outside the Roman one reads the same lesson at a
different moment of the year.} The Milan de Sirturis Ambrosian Missal of 1712
carries \latin{flecto genua} once, in a window headed \latin{sabbato in albis
missa pro baptizatis}, and the pericope there runs from \latin{Fratres obsecro
vos ne deficiatis} to \latin{in omnes generationes seculi seculorum amen} ---
\textbf{Eph.~3:13--21 entire, with the same liturgical incipit adaptation the
Roman book makes}, read as the baptismal lesson of Easter Saturday with the
neophyte chant \latin{quasi modo geniti infantes} in the same stream.

\textbf{Wilson and Schuster give the offset two different mechanisms}:
Wilson's, a difference of numeration base between \latin{post Pentecosten} and
\latin{post Octavas Pentecostes}; Schuster's, the interruption of the
post-Pentecost count by cycles reckoned from St~Lawrence, St Cyprian and
St~Michael. \textbf{No book collated for this Sunday decides between them.} Schuster's own title for this Sunday, \latin{Prima post
natale Sancti Cypriani}, is the hinge where the Lawrence count ends and the
Cyprian count begins --- \textbf{a title he prints at the head of his chapter
without naming the book he takes it from}, and he says elsewhere in the same
series that the count of these feast-relative Sundays was not uniform. The 1962
book prints only \latin{DOMINICA DECIMA SEXTA post Pentecosten}.

This Sunday also opens the temporal cycle's course of Ephesians, and the single
break in it falls at the Eighteenth Sunday, which Schuster judges to have been
originally aliturgical. \textbf{Schuster says the course runs until the
Twenty-third Sunday; the 1962 book reads Ephesians on the Sixteenth,
Seventeenth, Nineteenth, Twentieth and Twenty-first and Philippians thereafter},
so his extent holds of the captivity Epistles as a group, or of a differently
numbered series, and of nothing narrower.

\textbf{Three witnesses name stations for the days around this Sunday and none
for this one.} Schuster names one for the Seventeenth Sunday, one for the
Nineteenth and one for each of the three September Ember Days; Wilson's Appendix
names the same September stations in the books' own words --- \latin{ad S.
Petrum}, \latin{ad S. Mariam}, \latin{ad Apostolos}; and Gerbert's own p.~178
heads the September Ember Sunday \latin{Die Dom. Ad S. PETRUM} while his p.~184,
carrying this formulary's orations, has no station of any kind.

\subsection*{C2. Only the Introit's centonisation was made for this Mass, and the Alleluia has the weakest claim of the five to its place}

\textbf{Four of the five chants are printed elsewhere in the same 1962 book.}
The Gradual serves the Third, Fourth, Fifth and Sixth Sundays after the
Epiphany, and its own first verse returns as the Alleluia verse of the
Eighteenth Sunday two Sundays later. The Offertory has one other appointment,
the Friday after Lent~II. The Communion has two, the Thursday after Lent~IV and
a votive Mass \latin{pro quacumque necessitate}, this last standing in the 1962
book's own text layer, which the registry does not reach. And the Alleluia's
verse, Ps.~97:1, is the Introit antiphon of the Fourth Sunday after Easter and
stands as the psalm verse at four Introits besides --- Easter Thursday, the
Assumption, the Maternity of the Blessed Virgin Mary and the Octave of the
Nativity --- five further appointments spread across three seasons
(Paschaltide, Christmastide and the time after Pentecost) and three
sections of the book, where four of the Gradual's five are a block of
consecutive Epiphany Sundays and the fifth is the Alleluia verse of the
Eighteenth Sunday, two Sundays after this one. \textbf{The Introit is the one
exception}: no other formulary in the book takes its centonisation of Ps.~85:3
and 5 with v.~1, though the psalm itself is drawn on at the Fifteenth Sunday, at
the Friday of Lent~III, at the Ember Friday of Lent~I and at the Most Holy Name
of Jesus.

Those are measurements against the calendar registry named in the References,
save the votive Mass, which the registry does not carry; Schuster asserts three
of the appointments independently and the registry adds the others.

\textbf{The Alleluia is the weakest of the five.} At AMS~188, as the
gregorien.info database indexes Hesbert's six oldest Roman graduals, only three
of the six --- Mont-Blandin, Corbie and Senlis --- carry the formulary as a set,
and those three disagree about the Gradual: \latin{Timebunt gentes} at Corbie and
Senlis, \latin{Bonum est confidere} and \latin{Misit Dominus verbum suum} at
Mont-Blandin. Monza has a Gradual at that formulary and nothing else,
\latin{Liberasti nos}, which is a fourth answer to the same question;
Compiègne and Rheinau have no entry at it at all, and their indexes are not
truncated. The only alleluia any of the six records there is \latin{Laudate
Dominum omnes gentes}, in Senlis --- which is also what the Ottobonianus margin
gives, by a wholly independent route --- and \latin{Cantate Domino canticum
novum} has no AMS~188 entry anywhere in them. Wilson brackets the Ottobonianus alleluia
cues as later-hand additions; what those brackets assert about the manuscript is
not established, and the database's result stands without them.

\subsection*{C3. What the 1962 book prints around this formulary it settles by a general rule and not at this Sunday}

Three of the things a reader meets on these two pages --- the rank, the
\latin{Credo} and the Preface of the Most Holy Trinity --- are the book's
standing dispositions for a class of days, and the fourth, the single oration
where an older book printed three, is the one place where the day's own
appointment visibly changed.

\textbf{The formulary appoints one Collect, one Secret and one
Postcommunion}, where the 1862 Pustet at the same Sunday printed \latin{Secunda
Oratio. A cunctis nos.} and \latin{Tertia ad libitum.} with their matching
\latin{Alia Secreta} and \latin{Alia Postcommunio}. Under the 1862 book those
followed from its \latin{Rubricae generales}, tit.~II and \latin{De Orationibus}
nn.~2 and 11, which prescribed \latin{A cunctis} as the second oration from the
octave of Pentecost to Advent; under the 1960 Code the day takes one, by 1962
RG n.~434~b, \latin{in dominicis II classis, nulla alia admittitur oratio,
praeter commemorationem festi II classis}. \textbf{The step between them is the
general decree of the Sacred Congregation of Rites of 23~March 1955}, \work{De
rubricis ad simpliciorem formam redigendis}, at \work{Acta Apostolicae Sedis}
47 (1955) 218--224, binding \latin{kalendis Ianuariis anni 1956}. It does the
work twice over: Tit.~II n.~1, \latin{Gradus et ritus semiduplex supprimitur},
removes the rank on which the 1862 three-oration rule for this day depended, and
Tit.~V~a) n.~1, \textbf{\latin{Orationes pro diversitate temporum assignatae
abolentur}}, abolishes the seasonally assigned orations of which the
\latin{A cunctis} is one. \textbf{The books this Sunday was celebrated from
print the decree's own phrase inside this formulary}: the 1920 Vatican typical
edition and the pre-1955 Benziger both set \latin{Orationes pro diversitate
Temporum assignatae, ut supra} immediately after the Collect, with
\latin{Aliae Secretae} and \latin{Aliae Postcommuniones} matching. The 1960 Code
inherited that change; its own contribution here was the reclassification from
semiduplex, and then duplex, to second class.

\latin{Credo.} restates a general rule rather than a decision of this Sunday:
1962 RG n.~475~a orders the symbol \latin{in qualibet dominica}, the 1862
book ordered it \latin{in omnibus Dominicis per annum}, and the 1955 decree
narrowed the rule between them --- Tit.~V~b) n.~7, \latin{Credo dicitur
dumtaxat in dominicis et festis I classis} --- and left this Sunday inside it.
\textbf{The Preface of the Most Holy Trinity is likewise not proper to this
Mass}: 1962 RG n.~494~b gives it \latin{tamquam de Tempore \dots\ in omnibus
dominicis II classis extra tempus natalicium et paschale}, and the 1862 book
carries the same rule at its \latin{De Praefatione} n.~5 while printing no
preface direction inside this Sunday's formulary. \textbf{The printing change is
older than 1962}: \latin{Præfatio de \Sanctissima{} Trinitate.} already stands inside this
formulary in the 1920 typical edition and in the Benziger beside it, so it
arrived with the Pius X typical edition. The 1759 attribution of that Sunday
preface to Clement~XIII is the \work{Catholic Encyclopedia}'s report at its
article ``Preface''; the article names no \work{Acta} or \work{Decreta} locus,
and Fortescue's own \work{The Mass} carries neither the rule nor the date.

\textbf{Nothing printed around the ten appointed texts was chosen for this
Sunday}, and the one element of the frame that did change here is the one whose
removal a decree outside the Missal records.

\subsection*{C4. The published English reaches this Mass's doctrine at one element, and stops short at the two words its strongest arguments turn on}

\textbf{The Douay follows the psalter and not the chant.} It renders
\latin{gloria} at the Gradual and \latin{ad adjuvandum me} at the Offertory, so
it answers neither \latin{maiestáte} nor \latin{in auxílium meum réspice} ---
the two Latin words on which the strongest cross-element arguments of this
formulary turn --- and those arguments have to be made about the Latin, at the
element, and cannot be shown to an English reader in the appointed words at all.
\textbf{At the Gradual the argument survives the gap even though the word does
not}: Bellarmine's 1866 English states both comings at the verse itself, and
where his lemma keeps the psalter's noun his exposition reaches the chant's. At
the Offertory nothing answers, his English there being the Douay's.

\textbf{The Communion is the exception in both directions}: it departs from the
psalter at no word, and therefore its English says what its Latin says. The
Douay of Ps.~70:16 reads \emph{I will be mindful of thy justice alone}, which is
the antiphon's Latin word for word. \textbf{So the one place in this Mass where
its doctrine of grace stands in English on the appointed words is a chant and
not a prayer} --- and Augustine's exposition of Ps.~70:16--18 is an exposition of
the sung text and of nothing else.

\textbf{The three orations reach English by another route and lose more of the
Latin.} Their witness is a free devotional rendering of 1861 and not a construe:
\latin{prævéniat} arrives in the older sense of \emph{prevent}, which an English
reader now hears as its opposite; \latin{pérfice miserátus in nobis}, the verb
the Secret gives God and the clause that governs \latin{mereámur}, does not
survive into it at all; and the Postcommunion's \latin{córporum \dots\ præsens
páriter et futúrum \dots\ auxílium} becomes a paraphrase that keeps the two
times and loses the noun they belong to. \textbf{What the orations lose in that
rendering Guéranger supplies in his own}: he prints an English of the Collect
beside his exposition of it and reads the Secret, the Communion and the
Postcommunion in English at their own places, so the three orations and the
Communion antiphon all reach an English reader by a route the missal itself does
not print. \textbf{The
doctrine the three prayers say in Latin is sayable to an English reader on the
appointed words at the Communion and nowhere else in this Mass}, and the two
decisive chant words have no English witness here at all.

\clearpage
